Ce chapitre présente l'analyse et la conception détaillée de notre plateforme de gestion de projets.

\section{Besoins Fonctionnels}
\begin{itemize}
  \item Gestion des projets : création, suivi, types (Scrum/Itératif), états et budgets.
  \item Phases et jalons : planification, exécution, maîtrise, clôture avec milestones.
  \item Gestion agile : sprints, user stories, tâches avec priorités, états et affectations.
  \item Backlog et cahier des charges : gestion versionnée et documents contractuels.
  \item Ressources et affectations : rôles projet et composition d'équipe.
  \item Livrables et validations : workflow de validation et commentaires.
  \item Changements : demandes, priorisation, approbation/rejet et implémentation.
  \item Tableau de bord et analytics : KPI, vélocité, burndown, répartition des tâches.
  \item Notifications : événements clés en temps réel.
\end{itemize}

\section{Acteurs}
\begin{itemize}
  \item \textbf{Administrateur} : gouvernance globale, gestion des utilisateurs et paramètres.
  \item \textbf{Chef de projet / Product Owner} : création et pilotage des projets, sprints, backlog.
  \item \textbf{Scrum Master} : facilitation agile, suivi d'avancement, gestion des obstacles.
  \item \textbf{Développeur/Testeur} : exécution des tâches, mise à jour des statuts, commentaires.
  \item \textbf{Direction/Client} : validation des livrables, revue des KPI et décisions.
\end{itemize}

\section{Contraintes et Règles}
\begin{itemize}
  \item Cohérence des dates (phase/sprint/tâche) et états synchronisés.
  \item Permissions basées sur les rôles et la participation au projet.
  \item Traçabilité des changements d'état, validations et commentaires.
  \item Performance des listes paginées et filtrage.
\end{itemize}

\section{Modèles et Diagrammes}
\begin{itemize}
  \item Diagramme de cas d'utilisation : gestion de projet, sprint, backlog, livrable, validation.
  \item Diagramme de classes : entités \texttt{Projet, Phase, Sprint, UserStory, Tache, Ressource, Affectation, Livrable, Validation, Changement, Notification} et relations.
  \item Diagrammes de séquence : planification de sprint, création de tâche, validation de livrable.
\end{itemize}

\section{Règles Métier Illustratives}
\begin{itemize}
  \item Une tâche \texttt{critique=true} doit être visible en priorité et signalée dans les KPI.
  \item Un sprint ne peut être clôturé que si toutes les user stories sont terminées ou replanifiées.
  \item Un livrable ne peut passer à “Approuvé” qu'après validation par un rôle habilité.
  \item Une demande de changement “Approuvée” doit créer un suivi d'implémentation.
\end{itemize} 